\documentclass[10pt]{article}
\usepackage[T1]{fontenc}
\usepackage{newtxtext}
\usepackage{newtxmath}
\usepackage[letterpaper,top=0.68in,bottom=0.70in,left=0.72in,right=0.72in,
  headheight=12pt,headsep=7pt,footskip=20pt]{geometry}
\usepackage{amsmath,mathtools,bm}
\usepackage{booktabs,tabularx,array}
\usepackage{enumitem}
\usepackage{xcolor}
\usepackage{titlesec}
\usepackage{caption}
\usepackage{fancyhdr}
\usepackage{hyperref}
\usepackage{url}

% Self-contained compact OR-Bench layout.
\definecolor{orblue}{RGB}{24,67,108}
\definecolor{orgray}{RGB}{88,94,101}
\definecolor{orlight}{RGB}{246,247,248}
\hypersetup{
  colorlinks=true,
  linkcolor=orblue,
  citecolor=orblue,
  urlcolor=orblue,
  pdfborder={0 0 0}
}
\setlength{\parindent}{1em}
\setlength{\parskip}{0pt}
\setlength{\tabcolsep}{4.5pt}
\renewcommand{\arraystretch}{1.08}
\setlist[itemize]{leftmargin=1.35em,itemsep=1pt,topsep=3pt,parsep=0pt}
\setlist[enumerate]{leftmargin=1.45em,itemsep=1pt,topsep=3pt,parsep=0pt}
\captionsetup{font=small,labelfont=bf,labelsep=period,skip=4pt,hypcap=false}
\titleformat{\section}
  {\bfseries\fontsize{11.6}{13.2}\selectfont}
  {\thesection.}{0.45em}{}
\titleformat{\subsection}
  {\bfseries\fontsize{10.2}{11.8}\selectfont}
  {\thesubsection.}{0.45em}{}
\titlespacing*{\section}{0pt}{10pt plus 2pt minus 2pt}{3pt}
\titlespacing*{\subsection}{0pt}{7pt plus 1pt minus 1pt}{2pt}
\makeatletter
\newcommand{\obshorttitle}{OR-Bench Candidate}
\newcommand{\shorttitle}[1]{\renewcommand{\obshorttitle}{#1}}
\AtBeginDocument{\hypersetup{pdftitle={\@title},pdfauthor={\@author}}}
\renewcommand{\maketitle}{%
  \thispagestyle{plain}%
  \begin{center}
    {\fontsize{17}{19.5}\selectfont\bfseries\@title\par}
    \vspace{4pt}
    {\small\@author\par}
    \ifx\@date\@empty\else\vspace{1pt}{\footnotesize\color{orgray}\@date\par}\fi
  \end{center}
  \vspace{4pt}\hrule height 0.55pt\vspace{7pt}
}
\makeatother
\pagestyle{fancy}
\fancyhf{}
\fancyhead[L]{\scriptsize\color{orgray}\obshorttitle}
\fancyhead[R]{\scriptsize\color{orgray}OR-Bench candidate problem}
\fancyfoot[C]{\scriptsize\thepage}
\renewcommand{\headrulewidth}{0.35pt}
\renewcommand{\footrulewidth}{0pt}
\fancypagestyle{plain}{
  \fancyhf{}
  \fancyfoot[C]{\scriptsize\thepage}
  \renewcommand{\headrulewidth}{0pt}
}
\renewenvironment{abstract}{%
  \noindent\begin{minipage}{\linewidth}\small
  \textbf{Abstract.}\hspace{0.35em}%
}{%
  \end{minipage}\par\vspace{5pt}
}
\newcommand{\keywords}[1]{%
  \noindent{\small\textbf{Keywords.}\hspace{0.35em}#1}\par\vspace{5pt}
}
\newenvironment{businessbrief}{%
  \par\vspace{4pt}\noindent
  \begin{minipage}{\linewidth}
  \hrule height 0.4pt\vspace{4pt}
  \small\textbf{Business-level description.}\hspace{0.35em}\itshape
}{%
  \par\vspace{4pt}\hrule height 0.4pt
  \end{minipage}\par\vspace{4pt}
}
\newcommand{\evidencebox}[1]{%
  \par\vspace{3pt}\noindent
  \colorbox{orlight}{%
    \parbox{\dimexpr\linewidth-2\fboxsep\relax}{\small\textbf{Failure certificate.}\hspace{0.35em}#1}%
  }\par\vspace{4pt}
}
\newcommand{\metainfo}[1]{%
  \noindent{\footnotesize\color{orgray}#1}\par\vspace{5pt}
}
\newcolumntype{Y}{>{\raggedright\arraybackslash}X}
\newcolumntype{R}{>{\raggedleft\arraybackslash}X}
\AtBeginDocument{%
  \setlength{\abovedisplayskip}{6pt plus 2pt minus 2pt}
  \setlength{\belowdisplayskip}{6pt plus 2pt minus 2pt}
  \setlength{\abovedisplayshortskip}{3pt plus 1pt}
  \setlength{\belowdisplayshortskip}{4pt plus 1pt}
  \setlength{\jot}{2pt}
}
\providecommand{\doi}[1]{\href{https://doi.org/#1}{doi:#1}}

\title{Cyclic Freight Service Design with Cross-Dock Handling}
\author{Zhengzhong You}
\date{OR-Bench candidate problem \textbullet{} August 2026}
\shorttitle{Cyclic Cross-Dock Service Network Design}

\begin{document}
\maketitle

\begin{abstract}
Scheduled service network design must synchronize two flows: freight follows time-feasible connections, while reusable vehicles cover the weekly timetable with a finite fleet. This candidate isolates their interaction in an eight-leg periodic network with two tractors and one-slot cross-dock handling. The exact optimum is a direct out-and-back rotation costing 15. A generated GPT-5.4 program declares the feasible instance infeasible after duplicating destination demand. Two independent GPT-5-mini programs instead prove 11.5 with a design that needs three staggered tractors. Exact enumeration certifies both errors.
\end{abstract}
\keywords{service network design; cross-docking; periodic time-space network; fleet balance; mixed-integer linear programming}
\metainfo{Form domain: Supply chain and logistics \textbullet{} Model: MILP \textbullet{} Exact oracle: all $2^8$ service-leg designs}

\section{Business Brief}

\begin{businessbrief}
A freight carrier publishes the same four-slot line-haul plan every week. Each operated leg moves one owned tractor; unlisted deadheads and rentals are unavailable, but tractors may wait freely at terminals. At most two tractors cover the repeating timetable. One indivisible load is ready at O in slot 0 and is due at D by slot 3. At hub H, unloading, sorting, and reloading occupy one complete slot: a load arriving in slot 1 is ready only in slot 2. Choose the weekly service legs and load path that minimize operating cost.
\end{businessbrief}

The request is a compact scheduled service network design problem: vehicle conservation alone does not certify that the load can execute the same timetable.

\section{Periodic Network and Instance}

Let $N=\{O,H,D\}\times\{0,1,2,3\}$ be periodic time-space nodes, $A$ the listed tractor legs, and $B$ the zero-cost one-slot waiting arcs. A return leg enters its stated node in the next weekly cycle. Each service leg has frequency at most one.

\begin{center}
\captionof{table}{Candidate service legs. Return legs are cargo-ineligible.}
\begin{tabularx}{0.96\linewidth}{l Y Y r l@{}}
\toprule
Leg & Tail & Head & Cost & Cargo \\
\midrule
$F1$ & $O0$ & $H1$ & 2 & yes \\
$R1$ & $H1$ & $O0$ next cycle & 2 & no \\
$F2$ & $H1$ & $D2$ & 2 & yes \\
$R2$ & $D2$ & $H1$ next cycle & 2 & no \\
$F3$ & $H2$ & $D3$ & 10 & yes \\
$R3$ & $D3$ & $H2$ next cycle & 10 & no \\
$F4$ & $O0$ & $D3$ & 7.5 & yes \\
$R4$ & $D3$ & $O0$ next cycle & 7.5 & no \\
\bottomrule
\end{tabularx}
\end{center}

The one-slot handling transition connects the freight state $H1$ to $H2$. Thus the correct feasible cargo paths are
\[
\mathcal R=\{r^{\rm dir}=(F4),\;r^{\rm late}=(F1,F3)\}.
\]
The tempting pair $(F1,F2)$ is excluded: $F1$ arrives at 1, handling finishes at 2, and $F2$ departs at 1.

\section{Exact MILP}

Let $y_a\in\{0,1\}$ indicate that service leg $a$ operates, $w_b\in\mathbb Z_+$ count tractors on waiting arc $b$, and $z_r\in\{0,1\}$ select cargo path $r\in\mathcal R$. Write $d_e$ for arc duration, $H=4$ for the weekly horizon, and $M=2$ for the fleet. The exact formulation is
\begin{align}
\min\quad &\sum_{a\in A}c_a y_a, \label{eq:obj}\\
\text{s.t.}\quad
&\sum_{r\in\mathcal R}z_r=1, \label{eq:path}\qquad\\[-9pt]
&z_r\le y_a &&\forall r\in\mathcal R,\ a\in r, \label{eq:link}\\
&\sum_{a\in A\cap\delta^+(n)}y_a+\sum_{b\in B\cap\delta^+(n)}w_b
=\sum_{a\in A\cap\delta^-(n)}y_a+\sum_{b\in B\cap\delta^-(n)}w_b
&&\forall n\in N, \label{eq:balance}\\
&\sum_{a\in A}d_a y_a+\sum_{b\in B}d_b w_b\le HM=8, \label{eq:fleet}\\
&y_a,z_r\in\{0,1\},\quad w_b\in\mathbb Z_+. \label{eq:domain}
\end{align}
Equation~\eqref{eq:balance} closes the periodic tractor circulation. Equation~\eqref{eq:fleet} is essential: total tractor-time per period cannot exceed two tractor-weeks. Path eligibility is generated only after release, deadline, and the connection rule $\operatorname{dep}(b)\ge\operatorname{arr}(a)+h_H$ through H, where $h_H=1$.

\section{Exact Outcome and Diagnostic Ablations}

The direct path opens $F4,R4$ and costs 15. The legal transfer path is coverable by $F1$, one hub wait, $F3$, and $R4$, but costs 19.5. Thus the unique exact optimum is direct. Five exhaustive cases localize the two principal modeling layers.

\begin{center}
\captionof{table}{Correct model and four diagnostic ablations.}
\begin{tabularx}{0.96\linewidth}{Y l r}
\toprule
Model & Selected cargo/service legs & Objective \\
\midrule
Correct handling, circulation, fleet & $F4;\ F4,R4$ (1 tractor) & \textbf{15} \\
Omit handling only & $F1,F2;\ F1,F2,R1,R2$ (2) & 8 \\
Omit fleet limit only & $F4;\ F4,R1,R2$ (3) & 11.5 \\
Omit all asset constraints & $F4;\ F4$ & 7.5 \\
Omit handling and all assets & $F1,F2;\ F1,F2$ & 4 \\
\bottomrule
\end{tabularx}
\end{center}

\evidencebox{A GPT-5.4 (2026-03-05, low reasoning) program returns INFEASIBLE. Its shipment balance assigns supply $+1$ at $O0$ but demand $-1$ at both $D2$ and $D3$, so total demand is two for one load. The explicit design $F4,R4$ carries the load directly, uses one tractor, and costs 15; infeasibility is therefore false.}

A separate generated Gurobi model proves 11.5 with $F4,R1,R2$. Its modulo-week walk uses $F4$, three D-waits, $R2$, and $R1$: $3+3+3+3=12$ tractor-slots. Repeating its weekly departures needs $12/4=3$ staggered tractors, but only two exist. The underestimation is 3.5, or 23.33\% of the correct optimum.

Two independent GPT-5-mini (2025-08-07, high reasoning) runs on the final canonical prompt produced executable models with the same infeasible design and objective 11.5. One collapses the winding cycle into a named-tractor circulation; the other undercounts cross-week arcs. Contemporaneous high-reasoning runs of GPT-5.4, GPT-5.5, and GPT-5.6 Luna returned the correct objective 15. The package therefore records distinct, reproducible errors without claiming universal failure.

\section{Reproducibility and Sources}

The repository supplies the original synthetic instance, exact Gurobi model, independently generated shipment paths, a $2^8$ standard-library oracle, four ablations, the exact evaluation prompts, executable screened output, cleaned execution record, ground truth, and hashes. The statement and data are original CC BY 4.0 work; the solver is MIT licensed. No upstream data or code is copied. Public package: \url{https://github.com/Zhengzhong-You/OR-Bench}.

Periodic time-space SND and vehicle design balance follow the conceptual setting of Gao, Jin, and Diao~\cite{gao2022}. The one-slot loading, unloading, and transshipment semantics is motivated by Google's official Shipping Network Design documentation~\cite{google2024}.

\begin{thebibliography}{9}\small
\bibitem{gao2022}
A. Gao, X. Jin, and X. Diao. 2022.
An iterative backbone algorithm for service network design problems.
\emph{Processes} 10(7):1373. \doi{10.3390/pr10071373}.

\bibitem{google2024}
Google for Developers. 2024.
Shipping Network Design, Operations Research API.
\url{https://developers.google.com/optimization/service/shipping/network_design}.
\end{thebibliography}

\end{document}
